% Smoke test: apchem-notes.cls + every shared macro + the full graphics stack.
% If this compiles in all four variants, the infrastructure is sound.
\documentclass{apchem-notes}

\apcunit{0}
\apcunittitle{Infrastructure Smoke Test}
\apcdoctitle{Guided Notes}
\apcsubtitle{Exercises every macro once --- not real content}

\begin{document}
\apctitleblock

\blocklesson{Toggle and macro check}{Smoke --- 90 min}

\begin{objectives}
  \item Compile with \texttt{mhchem}, \texttt{siunitx}, \texttt{chemfig},
        \texttt{pgfplots}, \texttt{tikz}, \texttt{tcolorbox}.
  \item Render identically-sized blanks in student and key builds.
\end{objectives}

\reading{\S 7.12}{368--369}

\begin{warmup}[3]
  \item Write the formula for sulfuric acid: \answerblank[3cm]{\ce{H2SO4}}
  \item How many significant figures in \num{0.00340}? \answerblank[2cm]{3}
  \item State Avogadro's number: \answerblank[4cm]{\SI{6.022e23}{\per\mole}}
\end{warmup}

\segment{INSTRUCTION A}{25}{Chemistry rendering}

\notesection{Equations and units}{1.1}
\practice{5}

A balanced equation with states and an equilibrium arrow:
\[ \ce{N2(g) + 3 H2(g) <=> 2 NH3(g)} \qquad
   \Delta H = \SI{-92.2}{\kilo\joule\per\mole} \]

Redox half-reaction: \ce{Cu^2+ + 2 e- -> Cu}.
Molarity: \SI{0.150}{\mol\per\liter}. Range: \SIrange{25}{75}{\celsius}.

\term{Key vocabulary} renders in accent color. Fill in:
the mole is defined via \answerblank[6.4cm]{exactly $6.02214076\times10^{23}$ entities}.

\begin{workedexample}[Worked example: average atomic mass]
  Copper is 69.15\% \ce{^{63}Cu} (\SI{62.93}{u}) and 30.85\% \ce{^{65}Cu}
  (\SI{64.93}{u}).
  \[ 0.6915(62.93) + 0.3085(64.93) = \SI{63.55}{u} \]
\end{workedexample}

\begin{apctrap}
  Students average 62.93 and 64.93 to get 63.93, ignoring abundance weighting.
\end{apctrap}

\begin{apcnote}
  The equations sheet gives $E = h\nu$; students must supply $\nu = c/\lambda$.
\end{apcnote}

\segment{GUIDED PRACTICE}{15}{Structures and diagrams}

\notesub{A Lewis-adjacent structure via chemfig}

\chemfig{H-[:30]O-[:-30]H} \hspace{2em}
\chemfig{O=C=O}

\notesub{A particulate diagram via TikZ}

\begin{center}
\begin{tikzpicture}[scale=0.9]
  \draw[apcaccent, thick] (0,0) rectangle (4,2.4);
  \foreach \x/\y in {0.6/0.5, 1.5/1.8, 2.4/0.9, 3.3/1.6, 1.0/1.2}
    \fill[apcaccent!60] (\x,\y) circle (0.16);
  \foreach \x/\y in {0.9/1.9, 2.0/0.4, 2.9/1.9, 3.5/0.6}
    \draw[apcans, thick] (\x,\y) circle (0.16);
  \node[font=\sffamily\footnotesize, below] at (2,-0.1)
    {filled = \ce{Na+}, open = \ce{Cl-}};
\end{tikzpicture}
\end{center}

\segment{INSTRUCTION B}{20}{Plots}

A Maxwell--Boltzmann sketch via pgfplots:

\begin{center}
\begin{tikzpicture}
\begin{axis}[width=9cm, height=5cm,
    xlabel={molecular speed}, ylabel={fraction},
    xtick=\empty, ytick=\empty,
    axis lines=left, domain=0:1200, samples=80]
  \addplot[seriesA] {x^2*exp(-x^2/80000)/1500};
  \addplot[seriesB] {x^2*exp(-x^2/180000)/2500};
\end{axis}
\end{tikzpicture}
\end{center}

\segment{APPLICATION}{20}{Free-form answer space}

Explain why the higher-temperature curve flattens and shifts right:
\answerlines[3]{Same total area (same $N$), but the distribution spreads over
more speeds; the peak drops and the most-probable speed increases with $T$.}

Show your work below:
\workspace[2.5cm]

\begin{exitticket}
  One question, collected on the way out:
  \ce{^{20}Ne} vs \ce{^{22}Ne} --- same element? Why?
  \answerlines[2]{Yes: identical proton count (10); they differ only in neutrons.}
\end{exitticket}

\begin{spiralstrip}{drawn from earlier units}
  \item (U0) Convert \SI{2.5}{\gram} \ce{H2O} to moles. \answerblank[3cm]{\SI{0.14}{\mole}}
  \item (U0) \lmsonly{This sentence appears only in the LMS build.}%
        \printonly{This sentence appears only in the print build.}
\end{spiralstrip}

\end{document}
